\documentclass[11pt,a4paper]{ctexart}
\usepackage[margin=2.2cm]{geometry}
\usepackage{xcolor}
\usepackage{enumitem}
\usepackage{amssymb}
\usepackage[normalem]{ulem}
\usepackage{titlesec}
\usepackage{hyperref}
\definecolor{navy}{HTML}{1A365D}
\definecolor{ink}{HTML}{1F2933}
\definecolor{muted}{HTML}{66788A}
\definecolor{soft}{HTML}{EEF4FA}
\definecolor{gold}{HTML}{8A5A12}
\hypersetup{colorlinks=true,linkcolor=navy,urlcolor=navy}
\setlist{nosep}
\titleformat{\section}{\Large\bfseries\color{navy}}{}{0pt}{}
\titleformat{\subsection}{\large\bfseries\color{navy}}{}{0pt}{}
\pagestyle{plain}

\title{\textbf{九下构词法专项（网传版）}\\[4pt]
\large 人教版九年级下册英语（网传电子版）· 合成 / 派生 / 转化 / 缩写 · 仅供教师参考}
\author{初中英语资料站}
\date{2026-11-12}
\begin{document}
\maketitle

\section*{使用说明}
\begin{itemize}[leftmargin=1.4em]
\item ⚠ 网传版本，非官方教材，学生请勿使用，仅供教师参考。
\item 本讲义对应九下 Grammar 附录第五板块：构词法。
\item 共 5 讲语法、6 级 48 题。
\item 在线版：\url{https://1901151458-prog.github.io/english-materials/g9x-wordform/}
\end{itemize}

\begin{center}
\fbox{\parbox{0.95\textwidth}{\small\textbf{⚠ 网传版本 · 非官方教材 · 仅供教师参考}\\
本册教材尚未正式出版（官方版预计 2027 年发布），本讲义依据网传电子版整理，板块划分与题目设计可能存在偏差。\textbf{学生请勿使用}；仅供教师备课参考。正式内容请以 2027 年官方教材为准。}}
\end{center}
\bigskip
\section{合成法}
由两个或两个以上的词合成一个新词。
\begin{itemize}[leftmargin=1.4em]
\item 复合名词：newspaper（名+名）、hotline（形+名）
\item 复合动词：highlight（形+动）、overcome（副+动）、sunbathe（名+动）
\item 复合形容词：hard-working（形+动-ing）、man-made（名+动-ed）、worldwide（名+形）
\item 复合副词：anyway、downstairs、everywhere
\item 复合数词：fifty-four、one-third
\end{itemize}

\section{派生法：前缀}
\begin{center}
\begin{tabular}{p{2.4cm}p{4.6cm}p{5.4cm}}
\hline
\textbf{前缀} & \textbf{含义} & \textbf{例}\\
\hline
un- & 不、非（反向或反义） & unlucky, unhappy, untidy, unusual\\
dis- & 不、反面（否定） & disagree, disappoint, dishonest, dislike\\
bi- & 双、二（数量） & bicycle, bilateral, bilingual, binoculars\\
inter- & 相互、交互、在一起 & international, internet, interview\\
re- & 再次、重新 & recycle, retell, reuse, rewrite\\
tele- & 远距离的 & telephone, telescope, television\\
\hline
\end{tabular}
\end{center}

\section{派生法：后缀}
\begin{center}
\begin{tabular}{p{4.0cm}p{8.0cm}}
\hline
\textbf{后缀} & \textbf{示例}\\
\hline
-er / -or（职业的人） & farmer, reporter, worker, actor, director, visitor\\
-ist（主义者／职业者） & artist, dentist, pianist, scientist, tourist\\
-ness（状态） & happiness, illness, kindness, sadness\\
-ion / -tion（行为、过程、结果） & action, education, graduation, invention\\
-ful（充满……的） & beautiful, careful, grateful, harmful, peaceful\\
-less（没有；无） & careless, helpless, homeless, weightless\\
-able（能……的） & comfortable, enjoyable, recyclable, unforgettable\\
-al / -y / -ing / -ed / -ly & formal, natural / windy, noisy / boring / bored / widely, friendly\\
\hline
\end{tabular}
\end{center}

\section{转化法}
一个单词由一种词类转换为另一种词类，转化后的意义往往与之前的意义联系密切。
\begin{itemize}[leftmargin=1.4em]
\item 动词 → 名词：cry v. 哭 → (have a) cry n.；look v. 看 → (have a) look n.
\item 名词 → 动词：book n. 书 → book v. 预定；water n. 水 → water v. 浇
\item 形容词 → 动词：slow adj. 慢的 → slow (down) v.；calm adj. 镇静的 → calm v. 使平静
\end{itemize}

\section{缩写和简写}
主要采取「截头」「去尾」或者「又截头又去尾」的方法生成新词。
\begin{itemize}[leftmargin=1.4em]
\item telephone → phone；laboratory → lab；examination → exam；influenza → flu；advertisement → ad
\item 首字母缩写：AI、BCE、CE、CPC、GDP、PE、PRC、TCM、UN
\end{itemize}


\newpage
\section{分级练习（48 题）}
\noindent 按难度分级排列，先做基础，再做成分与扩展，最后做综合题。

\subsection*{第 1 级 · 合成法}
\noindent\textcolor{gray}{两个或以上的词合成一个新词。}\par\smallskip
\begin{enumerate}[leftmargin=*,itemsep=4pt]
\item newspaper 的构成是 \_\_\_。
\begin{enumerate}[label=\Alph*.,leftmargin=2em,itemsep=1pt,topsep=2pt]
\item 名词+名词
\item 形容词+名词
\item 动词+名词
\item 副词+名词
\end{enumerate}
\item hotline 的构成是 \_\_\_。
\begin{enumerate}[label=\Alph*.,leftmargin=2em,itemsep=1pt,topsep=2pt]
\item 形容词+名词
\item 名词+名词
\item 副词+动词
\item 名词+动词
\end{enumerate}
\item hard-working 的构成是 \_\_\_。
\begin{enumerate}[label=\Alph*.,leftmargin=2em,itemsep=1pt,topsep=2pt]
\item 形容词+动词-ing 形式
\item 名词+动词-ed
\item 副词+动词
\item 数词+数词
\end{enumerate}
\item man-made 的构成是 \_\_\_。
\begin{enumerate}[label=\Alph*.,leftmargin=2em,itemsep=1pt,topsep=2pt]
\item 名词+动词-ed 形式
\item 形容词+名词
\item 动词+名词
\item 副词+形容词
\end{enumerate}
\item worldwide 的构成是 \_\_\_。
\begin{enumerate}[label=\Alph*.,leftmargin=2em,itemsep=1pt,topsep=2pt]
\item 名词+形容词
\item 形容词+名词
\item 动词+副词
\item 数词+名词
\end{enumerate}
\item downstairs 的构成是 \_\_\_。
\begin{enumerate}[label=\Alph*.,leftmargin=2em,itemsep=1pt,topsep=2pt]
\item 副词+名词
\item 名词+名词
\item 形容词+动词
\item 数词+数词
\end{enumerate}
\item fifty-four 的构成是 \_\_\_。
\begin{enumerate}[label=\Alph*.,leftmargin=2em,itemsep=1pt,topsep=2pt]
\item 数词+数词
\item 名词+数词
\item 形容词+数词
\item 副词+数词
\end{enumerate}
\item something 属于哪种构词法？
\begin{enumerate}[label=\Alph*.,leftmargin=2em,itemsep=1pt,topsep=2pt]
\item 合成法
\item 派生法
\item 转化法
\item 缩写
\end{enumerate}
\end{enumerate}

\subsection*{第 2 级 · 派生法 · 前缀}
\noindent\textcolor{gray}{un-、dis-、bi-、inter-、re-、tele-。}\par\smallskip
\begin{enumerate}[leftmargin=*,itemsep=4pt]
\item un- 表示 \_\_\_。
\begin{enumerate}[label=\Alph*.,leftmargin=2em,itemsep=1pt,topsep=2pt]
\item 不、非（反向或反义）
\item 双、二
\item 再次、重新
\item 相互、交互
\end{enumerate}
\item dis- 表示 \_\_\_。
\begin{enumerate}[label=\Alph*.,leftmargin=2em,itemsep=1pt,topsep=2pt]
\item 不、反面（否定）
\item 远距离
\item 双、二
\item 再次
\end{enumerate}
\item bi- 表示 \_\_\_。
\begin{enumerate}[label=\Alph*.,leftmargin=2em,itemsep=1pt,topsep=2pt]
\item 双、二（数量）
\item 相互
\item 再次
\item 远距离
\end{enumerate}
\item inter- 表示 \_\_\_。
\begin{enumerate}[label=\Alph*.,leftmargin=2em,itemsep=1pt,topsep=2pt]
\item 相互、交互、在一起
\item 再次
\item 不、非
\item 远距离
\end{enumerate}
\item re- 表示 \_\_\_。
\begin{enumerate}[label=\Alph*.,leftmargin=2em,itemsep=1pt,topsep=2pt]
\item 再次、重新
\item 双、二
\item 不、非
\item 相互
\end{enumerate}
\item tele- 表示 \_\_\_。
\begin{enumerate}[label=\Alph*.,leftmargin=2em,itemsep=1pt,topsep=2pt]
\item 远距离的
\item 再次
\item 双、二
\item 不、非
\end{enumerate}
\item unhappy 中的 un- 属于 \_\_\_。
\begin{enumerate}[label=\Alph*.,leftmargin=2em,itemsep=1pt,topsep=2pt]
\item 前缀
\item 后缀
\item 合成
\item 转化
\end{enumerate}
\item international 中的 inter- 表示 \_\_\_。
\begin{enumerate}[label=\Alph*.,leftmargin=2em,itemsep=1pt,topsep=2pt]
\item 相互、交互
\item 远距离
\item 再次
\item 否定
\end{enumerate}
\end{enumerate}

\subsection*{第 3 级 · 派生法 · 名词后缀}
\noindent\textcolor{gray}{-er/-or、-ist、-ness、-ion/-tion。}\par\smallskip
\begin{enumerate}[leftmargin=*,itemsep=4pt]
\item -er / -or 表示 \_\_\_。
\begin{enumerate}[label=\Alph*.,leftmargin=2em,itemsep=1pt,topsep=2pt]
\item 从事某种职业的人
\item 处于某种状态
\item 行为的过程或结果
\item 充满……的
\end{enumerate}
\item -ist 表示 \_\_\_。
\begin{enumerate}[label=\Alph*.,leftmargin=2em,itemsep=1pt,topsep=2pt]
\item 主义者、从事某种职业者
\item 处于某种状态
\item 能……的
\item 没有；无
\end{enumerate}
\item -ness 表示 \_\_\_。
\begin{enumerate}[label=\Alph*.,leftmargin=2em,itemsep=1pt,topsep=2pt]
\item 处于某种状态
\item 从事某种职业的人
\item 行为的过程
\item 充满……的
\end{enumerate}
\item -ion / -tion 表示 \_\_\_。
\begin{enumerate}[label=\Alph*.,leftmargin=2em,itemsep=1pt,topsep=2pt]
\item 行为；行为的过程或结果、状态
\item 职业的人
\item 状态
\item 没有；无
\end{enumerate}
\item 下列哪个词是名词后缀 -er 构成的？
\begin{enumerate}[label=\Alph*.,leftmargin=2em,itemsep=1pt,topsep=2pt]
\item teacher
\item beautiful
\item careless
\item dangerous
\end{enumerate}
\item happiness 中的 -ness 是 \_\_\_。
\begin{enumerate}[label=\Alph*.,leftmargin=2em,itemsep=1pt,topsep=2pt]
\item 名词后缀
\item 形容词后缀
\item 副词后缀
\item 前缀
\end{enumerate}
\item invention 中的 -tion 是 \_\_\_。
\begin{enumerate}[label=\Alph*.,leftmargin=2em,itemsep=1pt,topsep=2pt]
\item 名词后缀
\item 形容词后缀
\item 动词后缀
\item 副词后缀
\end{enumerate}
\item 下列哪个词是名词？
\begin{enumerate}[label=\Alph*.,leftmargin=2em,itemsep=1pt,topsep=2pt]
\item artist
\item artistic
\item artistically
\item artistically’s
\end{enumerate}
\end{enumerate}

\subsection*{第 4 级 · 派生法 · 形容词/副词后缀}
\noindent\textcolor{gray}{-ful、-y、-ing、-ed、-al、-able、-less、-ly。}\par\smallskip
\begin{enumerate}[leftmargin=*,itemsep=4pt]
\item -ful 表示 \_\_\_。
\begin{enumerate}[label=\Alph*.,leftmargin=2em,itemsep=1pt,topsep=2pt]
\item 充满……的；具有……的
\item 没有；无
\item 能……的
\item 处于某种状态
\end{enumerate}
\item -less 表示 \_\_\_。
\begin{enumerate}[label=\Alph*.,leftmargin=2em,itemsep=1pt,topsep=2pt]
\item 没有；无
\item 充满
\item 能……的
\item 职业的人
\end{enumerate}
\item -able 表示 \_\_\_。
\begin{enumerate}[label=\Alph*.,leftmargin=2em,itemsep=1pt,topsep=2pt]
\item 能……的；具有……性质的
\item 没有；无
\item 充满
\item 处于状态
\end{enumerate}
\item -y 是什么后缀？
\begin{enumerate}[label=\Alph*.,leftmargin=2em,itemsep=1pt,topsep=2pt]
\item 形容词后缀
\item 名词后缀
\item 副词后缀
\item 动词后缀
\end{enumerate}
\item -ing 作后缀时构成 \_\_\_。
\begin{enumerate}[label=\Alph*.,leftmargin=2em,itemsep=1pt,topsep=2pt]
\item 形容词
\item 名词
\item 副词
\item 动词
\end{enumerate}
\item -ed 作后缀时构成 \_\_\_。
\begin{enumerate}[label=\Alph*.,leftmargin=2em,itemsep=1pt,topsep=2pt]
\item 形容词
\item 动词
\item 名词
\item 副词
\end{enumerate}
\item -al 是什么后缀？
\begin{enumerate}[label=\Alph*.,leftmargin=2em,itemsep=1pt,topsep=2pt]
\item 形容词后缀
\item 名词后缀
\item 动词后缀
\item 副词后缀
\end{enumerate}
\item -ly 通常构成 \_\_\_。
\begin{enumerate}[label=\Alph*.,leftmargin=2em,itemsep=1pt,topsep=2pt]
\item 副词（也可构成形容词）
\item 只能构成副词
\item 只能构成名词
\item 只能构成动词
\end{enumerate}
\end{enumerate}

\subsection*{第 5 级 · 转化法}
\noindent\textcolor{gray}{词类转换：动词→名词、名词→动词、形容词→动词。}\par\smallskip
\begin{enumerate}[leftmargin=*,itemsep=4pt]
\item book n. 书 → book v. 预定 属于 \_\_\_。
\begin{enumerate}[label=\Alph*.,leftmargin=2em,itemsep=1pt,topsep=2pt]
\item 转化法
\item 合成法
\item 派生法
\item 缩写
\end{enumerate}
\item water n. 水 → water v. 浇 属于 \_\_\_。
\begin{enumerate}[label=\Alph*.,leftmargin=2em,itemsep=1pt,topsep=2pt]
\item 转化法
\item 合成法
\item 派生法
\item 缩写
\end{enumerate}
\item (have a) look n. 看 属于 \_\_\_。
\begin{enumerate}[label=\Alph*.,leftmargin=2em,itemsep=1pt,topsep=2pt]
\item 转化法
\item 合成法
\item 派生法
\item 缩写
\end{enumerate}
\item slow adj. 慢的 → slow (down) v. 缓 属于 \_\_\_。
\begin{enumerate}[label=\Alph*.,leftmargin=2em,itemsep=1pt,topsep=2pt]
\item 转化法
\item 合成法
\item 派生法
\item 缩写
\end{enumerate}
\item calm adj. 镇静的 → calm (yourself) v. 使平静 属于 \_\_\_。
\begin{enumerate}[label=\Alph*.,leftmargin=2em,itemsep=1pt,topsep=2pt]
\item 转化法
\item 合成法
\item 派生法
\item 缩写
\end{enumerate}
\item 转化法的定义是 \_\_\_。
\begin{enumerate}[label=\Alph*.,leftmargin=2em,itemsep=1pt,topsep=2pt]
\item 一个单词由一种词类转换为另一种词类
\item 两个词合成一个新词
\item 加词缀构成新词
\item 截头去尾构成新词
\end{enumerate}
\item cry v. 哭 → (have a) cry n. 属于 \_\_\_。
\begin{enumerate}[label=\Alph*.,leftmargin=2em,itemsep=1pt,topsep=2pt]
\item 转化法
\item 合成法
\item 派生法
\item 缩写
\end{enumerate}
\item 转化后的意义通常与原来 \_\_\_。
\begin{enumerate}[label=\Alph*.,leftmargin=2em,itemsep=1pt,topsep=2pt]
\item 联系密切
\item 完全无关
\item 相反
\item 相同
\end{enumerate}
\end{enumerate}

\subsection*{第 6 级 · 缩写与综合}
\noindent\textcolor{gray}{截头去尾的简写与首字母缩写词。}\par\smallskip
\begin{enumerate}[leftmargin=*,itemsep=4pt]
\item phone 是 \_\_\_ 的简写。
\begin{enumerate}[label=\Alph*.,leftmargin=2em,itemsep=1pt,topsep=2pt]
\item telephone
\item photograph
\item phonetic
\item telegraph
\end{enumerate}
\item lab 是 \_\_\_ 的简写。
\begin{enumerate}[label=\Alph*.,leftmargin=2em,itemsep=1pt,topsep=2pt]
\item laboratory
\item labyrinth
\item labour
\item label
\end{enumerate}
\item exam 是 \_\_\_ 的简写。
\begin{enumerate}[label=\Alph*.,leftmargin=2em,itemsep=1pt,topsep=2pt]
\item examination
\item example
\item examine
\item excellent
\end{enumerate}
\item flu 是 \_\_\_ 的简写。
\begin{enumerate}[label=\Alph*.,leftmargin=2em,itemsep=1pt,topsep=2pt]
\item influenza
\item fluid
\item flush
\item flute
\end{enumerate}
\item ad 是 \_\_\_ 的简写。
\begin{enumerate}[label=\Alph*.,leftmargin=2em,itemsep=1pt,topsep=2pt]
\item advertisement
\item addition
\item address
\item advice
\end{enumerate}
\item UN 的全称是 \_\_\_。
\begin{enumerate}[label=\Alph*.,leftmargin=2em,itemsep=1pt,topsep=2pt]
\item United Nations
\item United States
\item Universal Network
\item United Network
\end{enumerate}
\item GDP 的全称是 \_\_\_。
\begin{enumerate}[label=\Alph*.,leftmargin=2em,itemsep=1pt,topsep=2pt]
\item gross domestic product
\item general data process
\item global development plan
\item great domestic power
\end{enumerate}
\item 缩写和简写主要采取的方法是 \_\_\_。
\begin{enumerate}[label=\Alph*.,leftmargin=2em,itemsep=1pt,topsep=2pt]
\item 截头、去尾或截头又去尾
\item 加词缀
\item 两个词合成
\item 词类转换
\end{enumerate}
\end{enumerate}

\newpage
\section*{答案与解析}
\small
\begin{enumerate}[leftmargin=1.6em,itemsep=2pt]
\item \textbf{W1-1 A}：news（名词）+ paper（名词）；教材合成法例句。
\item \textbf{W1-2 A}：hot（形容词）+ line（名词）；教材表。
\item \textbf{W1-3 A}：hard（形容词）+ working（动词-ing）；教材表。
\item \textbf{W1-4 A}：man（名词）+ made（动词-ed）；教材表。
\item \textbf{W1-5 A}：world（名词）+ wide（形容词）；教材表。
\item \textbf{W1-6 A}：down（副词）+ stairs（名词）；教材表。
\item \textbf{W1-7 A}：fifty + four 都是数词；教材表。
\item \textbf{W1-8 A}：some + thing 合成；教材复合代词例。
\item \textbf{W2-1 A}：教材表6：un-（不、非，表示反向或反义）。
\item \textbf{W2-2 A}：教材表6：dis-（不、反面，表示否定）。
\item \textbf{W2-3 A}：教材表6：bi-（双、二、表示数量）。
\item \textbf{W2-4 A}：教材表6：inter-（相互、交互、在一起）。
\item \textbf{W2-5 A}：教材表6：re-（再次、重新）。
\item \textbf{W2-6 A}：教材表6：tele-（远距离的）。
\item \textbf{W2-7 A}：un- 加在词前，是前缀；教材表6例词。
\item \textbf{W2-8 A}：教材表6例词 international。
\item \textbf{W3-1 A}：教材表7：-er / -or（表示「从事某种职业的人」，名词后缀）。
\item \textbf{W3-2 A}：教材表7：-ist（表示「主义者，从事某种职业者」等）。
\item \textbf{W3-3 A}：教材表7：-ness（表示「处于某种状态」，名词后缀）。
\item \textbf{W3-4 A}：教材表7：-ion / -tion（表示「行为；行为的过程或结果、状态」等）。
\item \textbf{W3-5 A}：teach + -er；教材表7 例词 farmer、worker 等。
\item \textbf{W3-6 A}：-ness 是名词后缀；教材表7。
\item \textbf{W3-7 A}：-tion 是名词后缀；教材表7。
\item \textbf{W3-8 A}：art + -ist 构成名词；教材表7例词。
\item \textbf{W4-1 A}：教材表7：-ful（充满……的；具有……的，形容词后缀）。
\item \textbf{W4-2 A}：教材表7：-less（表示「没有；无」，形容词后缀）。
\item \textbf{W4-3 A}：教材表7：-able（表示「能……的；具有……性质的」）。
\item \textbf{W4-4 A}：教材表7：-y（表性质，形容词后缀）。
\item \textbf{W4-5 A}：教材表7：-ing（形容词后缀），如 boring、interesting。
\item \textbf{W4-6 A}：教材表7：-ed（形容词后缀），如 bored、excited。
\item \textbf{W4-7 A}：教材表7：-al（形容词后缀），如 formal、natural。
\item \textbf{W4-8 A}：教材表7：-ly（副词或形容词后缀），如 widely、friendly。
\item \textbf{W5-1 A}：教材转化法例句：名词转化为动词。
\item \textbf{W5-2 A}：名词转化为动词；教材例句。
\item \textbf{W5-3 A}：look v. 看 → (have a) look n.；教材例句。
\item \textbf{W5-4 A}：形容词转化为动词；教材例句。
\item \textbf{W5-5 A}：形容词转化为动词；教材例句。
\item \textbf{W5-6 A}：教材定义。
\item \textbf{W5-7 A}：动词转化为名词；教材例句。
\item \textbf{W5-8 A}：教材：单词转化后的意义往往与之前的意义联系密切。
\item \textbf{W6-1 A}：教材缩写例：telephone → phone。
\item \textbf{W6-2 A}：教材缩写例：laboratory → lab。
\item \textbf{W6-3 A}：教材缩写例：examination → exam。
\item \textbf{W6-4 A}：教材缩写例：influenza → flu。
\item \textbf{W6-5 A}：教材缩写例：advertisement → ad。
\item \textbf{W6-6 A}：教材首字母缩写词表。
\item \textbf{W6-7 A}：教材首字母缩写词表。
\item \textbf{W6-8 A}：教材：主要采取「截头」「去尾」或者「一截截头又去尾」的方法生成新词。
\end{enumerate}

\end{document}
